% =============================================================
%  Chapter 4 — Material-Aware Risk-Sensitive
%              Continual Hamiltonian Learning
%  Rewritten per professor's feedback:
%  - what changes about p in Stage 2 (new sensing channels)
%  - hierarchy table as organizing spine
%  - annotated p/q update with R_theta visually separate
%  - three loops as three computational jobs
%  - gradient subsection dedicated
%  - persistent-state table
%  - CVaR argued, not "proved to dominate"
% =============================================================

\chapter{Material-Aware Risk-Sensitive Continual Hamiltonian Learning}
\label{chap:material}

\providecommand{\Ham}{\mathcal{H}}
\providecommand{\Potential}{\mathcal{R}}
\providecommand{\PhaseState}{z}
\providecommand{\Pos}{q}
\providecommand{\Mom}{p}
\providecommand{\Mass}{\mathrm{M}}
\providecommand{\J}{\mathbb{J}}

% =======================================================================
\section{Why a Richer Learner Is Needed}
\label{sec:ch4:motivation}
% =======================================================================

Chapter~\ref{chap:grlsnam} established that a geometry-only Hamiltonian
learner already provides structured, sample-efficient, continually
correctable navigation.
But three architectural gaps remain that cannot be closed by tuning the
existing framework.

\emph{Gap 1 - no material-risk force channel.}
The potential $\Potential^{(1)}$ responds only to signed geometric
clearances $d_i$.
Two configurations with identical obstacle geometry but different terrain
risk receive identical forces: the learner is structurally blind to
material hazards.

\emph{Gap 2 - no tail-risk objective.}
The training loss $\mathcal{L}^{(1)}$ optimizes expected trajectory
quality.
It cannot concentrate gradient signal on the rare catastrophic
material-contact episodes that define real navigation risk.

\emph{Gap 3 - coefficient space saturated.}
The coefficient vector $(\beta,\lambda,\{\alpha_i\},\gamma)$ has no
slot for material-risk weights.
Adding a new force type requires a new energy term, a new coefficient,
a new gradient pathway, and a new update law none of which fit the
current architecture.

The response is \emph{not} to replace the geometry-only framework.
It is to \textbf{enlarge it} along four coupled axes simultaneously:
stored energy, coefficient space, objective, and adaptation timescale.
The chapter's organizing principle is that each enlargement is a
\emph{necessity}, not a design choice.
Adding one new energy term forces new structure in all four axes.

% =======================================================================
\section{Enriched Problem Setting}
\label{sec:ch4:setting}
% =======================================================================

The navigation task is identical to Chapter~\ref{chap:grlsnam}.
The enrichment lies entirely in the environment context:
\begin{equation}
  \mathcal{E}_\tau
  = \bigl\{
      \vec{c}_{\mathrm{target}},\;
      \{d_i\}_{i=1}^{C},\;
      \underbrace{\tilde{r}(\vec{c})}_{\text{soft risk}},\;
      \underbrace{\phi(\vec{c})}_{\text{hard hazard SDF}},\;
      \underbrace{P_t(\vec{q}) \in \mathbb{R}^{6\times P\times P}}_{\text{local risk patch}}
    \bigr\},
  \label{eq:ch4:context}
\end{equation}
where $\tilde{r}(\vec{c})\in[0,1]$ is a continuous soft-risk score,
$\phi(\vec{c})$ is the signed distance to hard hazard zones, and
$P_t(\vec{q})$ is a 6-channel local crop
$(\tilde{r},\partial_x\tilde{r},\partial_y\tilde{r},h_{\mathrm{mask}},\phi,\|\nabla\phi\|)$
re-centred at $\vec{q}$ at every integration step.

In this chapter $\tilde{r}$ and $\phi$ are oracle inputs (Setting~2):
true material maps are provided; inferring them from raw imagery is
the Stage~3 future extension (§\ref{subsec:ch7:stage3}). 

Figure~\ref{fig:teaser} summarizes our setting and motivation: a geometry-only policy (Stage~1) is hazard-blind, while the oracle material-aware learner (Stage~2) uses ground-truth material risk maps to improve navigation safety.
Perception-based inference of material risk from raw RGB--IR or hyperspectral observations is the Stage~3 future extension (§\ref{subsec:ch7:stage3}); it is not performed in this chapter.

\begin{figure*}[htbp]
    \centering
    \includegraphics[width=0.75\textwidth]{figs/teaser.png}
\caption{
\textbf{From geometry-only navigation to oracle material-aware structured navigation.}
Geometry-only planning is hazard-blind when terrain regions share similar geometric structure but differ in traversal risk. In the Stage~2 setting of this chapter, the learner is given oracle access to a soft material-risk field $\tilde r$ and a hard-hazard boundary field $\phi$, and incorporates them into a factored Hamiltonian flow together with geometry and dissipation.
The resulting policy bends away from risky terrain, excludes hard hazards, and improves safety while preserving smooth goal-directed motion.}
    \label{fig:teaser}
\end{figure*}

% =======================================================================
\section{Hierarchy of Enrichments: The Organizing Spine}
\label{sec:ch4:hierarchy}
% =======================================================================

Table~\ref{tab:ch4:hierarchy} is the structural map of this chapter.
Every subsequent section fills in one or more rows.
The table makes the coupling explicit: each new energy term forces
consequences in all four axes of force, coefficient, objective, update law.
Removing any link breaks the others.

\begin{table}[htbp]
\centering
\caption{%
  \textbf{Four coupled enrichments from Stage~1 to Stage~2.}
  Each row is one link in the enrichment chain.
  Reading left to right: a new energy term (col~1) forces a new force
  channel (col~2), a new coefficient (col~3), a new objective term
  (col~4), and a new update law (col~5).
}
\label{tab:ch4:hierarchy}
\small
\setlength{\tabcolsep}{3.8pt}
\renewcommand{\arraystretch}{1.4}
\resizebox{\textwidth}{!}{
\begin{tabular}{lllll}
\toprule
\textbf{Structured component} &
\textbf{Force channel} &
\textbf{Coeff.} &
\textbf{Objective term} &
\textbf{Update law} \\
\midrule
$H_{\mathrm{geom}}$ (Stage~1)
  & $F_{\mathrm{goal}}, F_{\mathrm{bar}}$
  & $\beta,\{\alpha_i\}$
  & $\mathcal{L}_{\mathrm{traj}}, \mathcal{L}_{\mathrm{clear}}$
  & Segment Gauss-Newton \\
\midrule
$H_{\mathrm{mat}}^{\mathrm{soft}} = \lambda_s\tilde{r}(\vec{c})$
  & $F_{\mathrm{soft}} = -\lambda_s\nabla\tilde{r}$
  & $\lambda_s$
  & $w_r\sum_t\tilde{r}(\vec{c}_t)\Delta s_t$
  & CVaR episode gradient \\
$H_{\mathrm{mat}}^{\mathrm{hard}} = \lambda_h b(\phi)$
  & $F_{\mathrm{hard}} = -\lambda_h\nabla b(\phi)$
  & $\lambda_h$
  & $w_h\sum_t\mathbf{1}[\phi<1\text{m}]$
  & CVaR episode gradient \\
$R_\theta$ (dissipative port)
  & $-\gamma\Mass^{-1}\Mom$ (separate)
  & $\gamma$
  & $w_\gamma|\gamma-\gamma_0|^2$
  & Friction alignment \\
$\mathcal{R}_{\nabla H}$ (field reg.)
  & (smoothness on $F_{\mathrm{tot}}$)
  & —
  & $w_c\mathcal{L}_{\mathrm{clear}}+w_m\mathcal{L}_{\mathrm{multi}}$
  & Curriculum passivity \\
\midrule
\textbf{Stage~2 total}
  & $F_{\mathrm{tot}}$: 5 channels
  & $(\beta,\lambda_s,\lambda_h,$\;$\{\alpha_i\},\gamma)$
  & $\mathcal{L}_{\mathrm{P2}}$: CVaR $+$ scaffold $+$ reg.
  & 3-loop hierarchy \\
\bottomrule
\end{tabular}}
\end{table}

% =======================================================================
\section{What Changes About \texorpdfstring{$p$}{p} in Stage~2}
\label{sec:ch4:p_role}
% =======================================================================

Chapter §\ref{sec:ch3:p_role} established that the
costate $\Mom$ simultaneously serves as geometric accumulator,
progress evaluator, force predictor, correction receiver, and port
amender.
In Stage~2, three of these roles are enriched.

\paragraph{New sensing channels flowing into $\Mom$.}
The costate update now receives two additional force contributions
(see \eqref{eq:ch4:p_update} below):
\begin{itemize}
  \item $-\tau\lambda_s\nabla_{\vec{c}}\tilde{r}$: pushes $\Mom$ away
        from high-risk terrain at \emph{every integration step},
        continuously accumulating a cotangent-space summary of the
        material-risk geometry encountered so far.
  \item $-\tau\lambda_h\nabla_{\vec{c}}b(\phi)$: injects a sharp
        repulsive impulse as $\phi\to 0$, encoding proximity to
        hard hazard boundaries.
\end{itemize}
These are new \textbf{sensing channels} (not just loss terms) that flow
into $\Mom$ as the trajectory unfolds.
In Stage~1, $\Mom$ accumulated only geometric obstacle information.
In Stage~2, it simultaneously accumulates geometric \emph{and}
material-risk information.

\paragraph{Richer energy balance and the passivity monitor.}
The stored energy $H_\theta^{(2)}$ now includes $H_{\mathrm{mat}}$,
so the energy increment $\Delta H_t = H_\theta^{(2)}(\PhaseState_{t+1})
- H_\theta^{(2)}(\PhaseState_t)$ reflects \emph{both} geometric and
material-risk dynamics.
When the material heads $(\lambda_s,\lambda_h)$ are first unfrozen
at the Phase~1→2 transition, $\Delta H_t$ spikes temporarily as the new
force channels inject energy into $\Mom$ before the dissipative port
$R_\theta$ has been re-calibrated.
The passivity surrogate $\mathcal{L}_{\mathrm{pass}} = \frac{1}{T}\sum_t
\mathrm{softplus}(\Delta H_t)$ detects this spike and delays curriculum
advancement until stability is restored.
$\Mom$ is therefore monitoring a richer energy balance in Stage~2.

\paragraph{The evaluator remains internal.}
In actor-critic RL, a separate critic network evaluates trajectories
\emph{after} execution, outside the dynamics.
In Stage~2, as in Stage~1, $\Mom$ performs evaluation \emph{inside}
the evolving trajectory: it senses material geometry, monitors energy
dissipation, predicts the next force balance via
$\dot{\vec{q}} = \Mass^{-1}\Mom$, and receives the CVaR gradient
signal through the rollout sensitivity recursion.
Adding the three loops of §\ref{sec:ch4:loops} does not add an
external evaluator but \textbf{enriches the internal one}. Figure~\ref{fig:ch4:p_stage2} summarizes this enrichment visually:
the new material-aware learner does not change the role of $\Mom$
qualitatively, but expands the information it accumulates and the
force balance it predicts.

\begin{figure}[t]
    \centering
    \includegraphics[width=0.96\textwidth]{figs/ch4_p_stage2.png}
    \caption{
    In the geometry-only learner, the cotangent variable $\Mom_t$ already acts as an internal evaluator and guide for local motion.
    In Stage~2, $\Mom_t$ is enriched by two additional sensing channels: a soft material-risk gradient and a hard-hazard boundary signal.
    These are integrated together with obstacle geometry, local clearance, dissipation, and progress monitoring into a cotangent-space summary used to predict the next force balance.
    The resulting half-step state $\Mom_{t+1/2}$ enables a corrected and smoother update of the configuration $\Pos_{t+1}$.
    The figure emphasizes that Stage~2 does not add an external evaluator; it enriches the internal one carried by $\Mom$.
    }
    \label{fig:ch4:p_stage2}
\end{figure}

\subsection{Phase-Space Update}
\label{subsec:ch4:dynamics}

The discrete dynamics read:
\begin{equation}
\boxed{
\begin{aligned}
  \Mom_{t+1}
  &= \Mom_t
  \underbrace{-\,\tau\nabla_{\Pos}H_{\mathrm{geom}}(\vec{q}_t,\Mom_t)
  }_{\text{conservative geom. force (Stage 1)}}
  \underbrace{-\,\tau\nabla_{\Pos}H_{\mathrm{mat}}(\vec{q}_t)
  }_{\text{new material channels}}
  \underbrace{-\,\tau R_\theta(\vec{q}_t,\Mom_t)
  }_{\substack{\text{dissipative port} \\ \text{SEPARATE from }H_\theta}}
  \underbrace{+\,\tau\mathbf{B}u_t,
  }_{\text{control port}} \\[4pt]
  \vec{q}_{t+1}
  &= \vec{q}_t + \tau\,\Mass^{-1}\Mom_{t+1},
\end{aligned}
}
\label{eq:ch4:p_update}
\end{equation}
where $\nabla_{\Pos}H_{\mathrm{mat}} = \lambda_s\nabla\tilde{r}
+ \lambda_h\nabla b(\phi)$ and
$R_\theta(\vec{q},\Mom) = \gamma\Mass^{-1}\Mom$ with $\gamma\geq 0$.

The visual separation in \eqref{eq:ch4:p_update} is deliberate:
the stored energy $H_\theta = H_{\mathrm{geom}} + H_{\mathrm{mat}}$
and the dissipative port $R_\theta$ are structurally distinct objects.
$H_\theta$ defines the conservative force topology;
$R_\theta$ removes excess kinetic energy.
Conflating them would entangle the stability analysis with the
task-objective design and make the gradient pathways of §\ref{sec:ch4:gradients}
opaque.

\paragraph{Energy balance.}
Along any trajectory of \eqref{eq:ch4:p_update}:
\begin{equation}
  \Delta H_t
  = -\tau\Mom_t^\top R_\theta(\vec{q}_t,\Mom_t)
    + \tau\Mom_t^\top\mathbf{B}u_t
  \;\leq\; \tau\Mom_t^\top\mathbf{B}u_t.
  \label{eq:ch4:energy_balance}
\end{equation}
The system is \textbf{passive} (stored energy non-increasing in the
absence of external input) when $R_\theta\succ 0$,
i.e., when $\gamma>0$.
This is monitored by $\mathcal{L}_{\mathrm{pass}}$.

\subsection{Observer and Progress Signals}
\label{subsec:ch4:observer}

The context update between stages is:
\begin{equation}
  y_{t+1} = \Psi(y_t, \vec{q}_{t+1};\, Z^{(1)}, P_{t+1}),
  \label{eq:ch4:observer}
\end{equation}
re-centring the obstacle extractor and risk patch at $\vec{q}_{t+1}$.
Two scalar progress signals are maintained per step:
\begin{equation}
  s_t^{\mathrm{seg}} = |\vec{q}_t - \vec{g}_k|
  \quad\text{(spatial progress)},
  \qquad
  s_t^{\mathrm{pass}} = \Delta H_t
  \quad\text{(energy dissipation / stability proxy)}.
  \label{eq:ch4:progress}
\end{equation}
$s_t^{\mathrm{seg}}$ feeds the curriculum conquest criterion;
$s_t^{\mathrm{pass}}$ feeds the passivity check.
Both are lightweight scalar signals with no stored trajectory.

% =======================================================================
\section{Factored Stored Energy and Force Channels}
\label{sec:ch4:energy}
% =======================================================================
Figure~\ref{fig:ch4:factored_hamiltonian} gives the chapter's
force-level interpretation of this decomposition: each added energy
term introduces a distinct force channel and therefore a distinct
behavioral effect on the trajectory.
The stored energy decomposes additively, with $R_\theta$ separate:

\begin{figure}[t]
    \centering
    \includegraphics[width=0.96\textwidth]{figs/ch4_factored_hamiltonian.png}
    \caption{
    \textbf{Factored Hamiltonian and force channels in Stage~2.}
    The stored energy decomposes into kinetic, geometry, and material components,
    \(
    H_\theta = H_{\mathrm{kin}} + H_{\mathrm{geom}} + H_{\mathrm{mat}},
    \)
    while the dissipative port
    \(
    R_\theta(\Pos,\Mom)=\gamma \Mass^{-1}\Mom
    \)
    remains structurally separate.
    The geometry term contributes goal-attractive and obstacle-repulsive forces, while the material term contributes soft-risk and hard-hazard repulsion.
    Together with the dissipative force, these channels define the total vector field that drives the phase-space rollout.
    The right panel illustrates the resulting trajectory-level effect: the robot is simultaneously attracted toward the goal, repelled from obstacles, deflected away from risky terrain, excluded from hard hazards, and smoothed by dissipation.
    }
    \label{fig:ch4:factored_hamiltonian}
\end{figure}

\begin{equation}
  H_\theta(\vec{q},\Mom)
  = \underbrace{\tfrac{1}{2}\Mom^\top\Mass^{-1}\Mom}_{H_{\mathrm{kin}}}
  + \underbrace{\beta D_\mathcal{G}(\vec{q},\vec{g})
    + \sum_j\alpha_j b_{\mathrm{IPC}}(d_j)
  }_{H_{\mathrm{geom}}\;\text{(Stage 1)}}
  + \underbrace{\lambda_s\tilde{r}(\vec{c})
    + \lambda_h b_{\mathrm{sp}}(\phi(\vec{c}))
  }_{H_{\mathrm{mat}}\;\text{(new in Stage 2)}},
  \label{eq:ch4:ham}
\end{equation}
with $b_{\mathrm{IPC}}$ the incremental potential contact barrier
and $b_{\mathrm{sp}}(s) = \mathrm{softplus}(-s/\sigma_\phi)$ the
hazard SDF barrier.

The five force channels and their coefficients are:
\begin{alignat}{3}
  F_{\mathrm{goal}}  &= -\beta\nabla D_\mathcal{G},
  &\quad& \text{coeff.\ } \beta,
  &\quad& \text{shapes goal-attracting basin;} \label{eq:ch4:F_goal}\\
  F_{\mathrm{bar}}   &= -\sum_j\alpha_j\nabla b_{\mathrm{IPC}}(d_j),
  &\quad& \text{coeff.\ } \{\alpha_j\},
  &\quad& \text{shapes obstacle-repulsive walls;} \label{eq:ch4:F_bar}\\
  F_{\mathrm{soft}}  &= -\lambda_s\nabla\tilde{r},
  &\quad& \text{coeff.\ } \lambda_s,
  &\quad& \text{deflects from soft-risk terrain;} \label{eq:ch4:F_soft}\\
  F_{\mathrm{hard}}  &= -\lambda_h\nabla b_{\mathrm{sp}}(\phi),
  &\quad& \text{coeff.\ } \lambda_h,
  &\quad& \text{repels from hard-hazard boundaries (steep barrier);} \label{eq:ch4:F_hard}\\
  -R_\theta          &= -\gamma\Mass^{-1}\Mom,
  &\quad& \text{coeff.\ } \gamma,
  &\quad& \text{removes kinetic energy (separate port).} \label{eq:ch4:R}
\end{alignat}

The total force on the momentum is:
$\Mom_{t+1} = \Mom_t + \tau(F_{\mathrm{goal}} + F_{\mathrm{bar}}
+ F_{\mathrm{soft}} + F_{\mathrm{hard}} - R_\theta)$,
which is \eqref{eq:ch4:p_update} written component-wise.

\begin{lemma}[Local risk-deflection under material-aware force channels]
\label{lem:ch4:risk_deflect}
Let $v_1, v_2 \in T_{\vec{q}}\mathcal{Q}$ be two feasible motion
directions satisfying
$\langle\nabla_{\vec{q}}H_{\mathrm{geom}},v_1\rangle =
 \langle\nabla_{\vec{q}}H_{\mathrm{geom}},v_2\rangle$
(equal geometric cost to first order).
If $\langle\nabla\tilde{r}(\vec{q}),v_1\rangle >
   \langle\nabla\tilde{r}(\vec{q}),v_2\rangle$,
then the Stage~2 momentum update~\eqref{eq:ch4:p_update} satisfies
\[
  \langle\Mom_{t+1}^{(\mathrm{S2})},v_1\rangle <
  \langle\Mom_{t+1}^{(\mathrm{S2})},v_2\rangle,
\]
while Stage~1 treats $v_1$ and $v_2$ equally.
Consequently, the Stage~2 configuration step
$\vec{q}_{t+1} = \vec{q}_t + \tau\Mass^{-1}\Mom_{t+1}$
is biased toward the lower-risk direction $v_2$.
\end{lemma}
\begin{proof}
The momentum difference between Stage~2 and Stage~1 is
$\Mom_{t+1}^{(\mathrm{S2})} - \Mom_{t+1}^{(\mathrm{S1})}
 = -\tau\lambda_s\nabla\tilde{r}(\vec{q}_t)
   -\tau\lambda_h\nabla b_{\mathrm{sp}}(\phi(\vec{q}_t))$
(from \eqref{eq:ch4:p_update}, \eqref{eq:ch4:F_soft}--\eqref{eq:ch4:F_hard}).
Taking the inner product with $v_1 - v_2$ gives
$-\tau\lambda_s\langle\nabla\tilde{r},\,v_1-v_2\rangle < 0$,
since $\lambda_s>0$ at Stage~2 convergence and
$\langle\nabla\tilde{r},v_1\rangle>\langle\nabla\tilde{r},v_2\rangle$.
Stage~1 carries no $\nabla\tilde{r}$ term, so it assigns equal
momentum along both directions.
Full details, including the $\lambda_h$ contribution, are in
Appendix~\ref{app:lem:risk_deflect}.
\end{proof}

% =======================================================================
\section{Objective Design: Terminal Correctness, Risk, and Tail Sensitivity}
\label{sec:ch4:objective}
% =======================================================================

\subsection{Episode Cost}
\label{subsec:ch4:J}

The episode cost accumulates over all integration steps:
\begin{equation}
  J(\theta)
  = w_g D_\mathcal{G}(\vec{q}_T,\vec{g})
  + w_\ell\sum_t\|\vec{q}_{t+1}-\vec{q}_t\|
  + w_r\sum_t\tilde{r}(\vec{c}_t)\Delta s_t
  + w_h\sum_t\mathbf{1}[\phi(\vec{c}_t)<1\,\text{m}].
  \label{eq:ch4:J}
\end{equation}
The four terms encode: terminal goal error, path length,
cumulative soft-risk exposure, and hard-hazard incursion count.

\paragraph{Terminal correctness.}
The first term $w_g D_\mathcal{G}(\vec{q}_T,\vec{g})$ provides
a proxy for field correctness without requiring ground-truth force
fields: if the terminal distance is small, the energy landscape was
shaping the trajectory toward the right target.
This is how GRL-SNAM validates that the learned field is
``correct enough'' without a ground-truth vector field to compare against.

\subsection{Why Tail-Risk? Expected Cost Is Insufficient}
\label{subsec:ch4:cvar_motivation}

The primary navigation loss is $\mathrm{CVaR}_{0.95}(J)$
rather than $\mathbb{E}[J]$.
The distinction matters: optimizing $\mathbb{E}[J]$ can leave a small
fraction of catastrophic material-contact episodes in the optimal policy.

\begin{proposition}[Expected cost allows tail failures]
\label{prop:ch4:cvar}
Let $J$ be the episode cost for a policy $\pi$ and suppose a fraction
$\epsilon=0.05$ of rollouts have cost $J_{\mathrm{cat}}\gg\mathbb{E}[J]$
(e.g., hard-hazard entries).
Then $\pi$ minimizing $\mathbb{E}[J]$ need not reduce
$J_{\mathrm{cat}}$: the optimizer trades off against the
$(1-\epsilon)$ majority of low-cost rollouts.
Formally,
$\mathrm{CVaR}_{1-\epsilon}(J) = \mathbb{E}[J\mid J\geq\mathrm{VaR}_{1-\epsilon}(J)]$
upper-bounds $\mathbb{E}[J]$ with equality only when $J$ is
deterministic, and concentrates gradient signal on the
catastrophic tail.
\end{proposition}

Optimizing $\mathrm{CVaR}_{0.95}(J)$ forces the learning to focus
on the worst 5\% of rollouts.
In the DFC2018 benchmark, these are precisely the episodes that
route through water bodies or railways, rare in the training
distribution but catastrophic in deployment.
An expected-cost surrogate treats these episodes as 5\% of the
gradient signal; CVaR makes them 100\%.

\begin{proposition}[CVaR descent implies continual risk-path improvement]
\label{prop:ch4:cvar_descent}
Let $\hat{\mathcal{L}}_\beta(\theta)=\widehat{\mathrm{CVaR}}_\beta(J(\theta))$
be the empirical CVaR over a batch $\mathcal{B}$.
If $\theta$ is updated by the subgradient step
$\theta_{k+1}=\theta_k - \eta_k\hat{g}_k$,
where $\hat{g}_k\in\partial_\theta\hat{\mathcal{L}}_\beta(\theta_k)$
is computed via the sensitivity recursion
\eqref{eq:ch4:dL_dlambda_s}--\eqref{eq:ch4:dL_dlambda_h},
then for $\eta_k$ sufficiently small:
\[
  \hat{\mathcal{L}}_\beta(\theta_{k+1})
  \leq \hat{\mathcal{L}}_\beta(\theta_k) + O(\eta_k^2).
\]
Moreover, under $\eta_k=\eta/\sqrt{k}$ with subgradient norm
$\|\hat{g}_k\|\leq G$ (bounded by
$\mathcal{L}_{\mathrm{pass}}\leq\varepsilon_{\mathrm{pass}}$),
the iterates converge to a \emph{stationary point} of
$\hat{\mathcal{L}}_\beta$ at rate
\[
  \min_{k\leq K}\|\hat{g}_k\| \leq \frac{C}{\sqrt{K}},
\]
where $C$ depends on $G$ and the initial objective gap.
\end{proposition}
\begin{proof}
$\mathrm{CVaR}_\beta$ is convex in the scalar $J$
(Rockafellar--Uryasev 2000), but the composition
$\theta\mapsto\mathrm{CVaR}_\beta(J(\theta))$ is generally nonconvex
because $J(\theta)$ is nonlinear in $\theta$ through the rollout dynamics.
The passivity constraint $\mathcal{L}_{\mathrm{pass}}\leq\varepsilon_{\mathrm{pass}}$
bounds $\|\nabla_{\vec{q}}H_\theta\|$, hence $\|\nabla_\theta J\|$,
hence $G$.
For a nonconvex Lipschitz objective with bounded subgradient, the
diminishing-step subgradient method converges to a stationary point
at rate $O(1/\sqrt{K})$.
This bound applies to Loop~2 (episodic CVaR gradient);
the two-timescale structure ensures Loop~1 is approximately converged
when Loop~2 updates $(\lambda_s,\lambda_h)$.
Full proof is in Appendix~\ref{app:prop:cvar_descent}.
\end{proof}

\subsection{Phase Objectives}
\label{subsec:ch4:phase_obj}

Phase~1 trains the geometry backbone:
\begin{equation}
  \mathcal{L}_{\mathrm{P1}}
  = w_t\mathcal{L}_{\mathrm{traj}}
  + w_v\mathcal{L}_{\mathrm{vel}}
  + w_\gamma|\gamma-\gamma_0|^2
  + w_c\mathcal{L}_{\mathrm{clear}}
  + w_m\mathcal{L}_{\mathrm{multi}}.
  \label{eq:ch4:LP1}
\end{equation}
Phase~2 adds the CVaR objective with a scaffold:
\begin{equation}
  \mathcal{L}_{\mathrm{P2}}
  = \mathrm{CVaR}_{0.95}(J)
  + \underbrace{0.3\,\mathcal{L}_{\mathrm{P1}}}_{\text{geometry scaffold}}
  + w_\lambda\underbrace{(\|\lambda_s\|^2+\|\lambda_h\|^2)}_{\text{material head reg.}}.
  \label{eq:ch4:LP2}
\end{equation}
The 0.3 scaffold weight retains geometry backbone stability when
$(\lambda_s,\lambda_h)$ are unfrozen: without it, the material
heads could overpower geometry forces before the risk terrain has
been sufficiently sampled.
The coefficient regularizer $w_\lambda$ prevents $\lambda_s$
and $\lambda_h$ from growing large before the training distribution
covers high-risk regions.

% =======================================================================
\section{Gradient-Field Learning Law}
\label{sec:ch4:gradients}
% =======================================================================

\begin{quote}\itshape
The policy is the integrated vector field of $H_\theta$.
Learning proceeds by reshaping that field, not by attaching losses
to final trajectories alone.
The gradient of any rollout-dependent loss must propagate through
the integrator and because $R_\theta$ is a separate port, its
gradient stream is decoupled from $H_\theta$'s.
\end{quote}

\subsection{Rollout Sensitivity Recursion}
\label{subsec:ch4:sensitivity}

Differentiating the update \eqref{eq:ch4:p_update} with respect to
parameters $\theta = (\beta,\lambda_s,\lambda_h,\{\alpha_j\},\gamma)$:

\begin{align}
  \frac{\partial\Mom_{t+1}}{\partial\theta}
  &= \frac{\partial\Mom_t}{\partial\theta}
  \underbrace{-\tau\frac{\partial\nabla_\Pos H_{\mathrm{geom}}}{\partial\theta}
  }_{\text{geom. gradient stream}}
  \underbrace{-\tau\frac{\partial\nabla_\Pos H_{\mathrm{mat}}}{\partial\theta}
  }_{\text{material gradient stream}}
  \underbrace{-\tau\frac{\partial R_\theta}{\partial\theta}
  }_{\substack{\text{dissipation stream}\\\text{(independent of }H_\theta\text{)}}},
  \label{eq:ch4:dp_dtheta} \\
  \frac{\partial\vec{q}_{t+1}}{\partial\theta}
  &= \frac{\partial\vec{q}_t}{\partial\theta}
  + \tau\Mass^{-1}\frac{\partial\Mom_{t+1}}{\partial\theta}.
  \label{eq:ch4:dq_dtheta}
\end{align}

The three additive streams in \eqref{eq:ch4:dp_dtheta} are structurally
independent because $H_\theta$ and $R_\theta$ are separate objects.
The dissipation stream $\partial R_\theta/\partial\gamma =
\Mass^{-1}\Mom_t$ is always available analytically; it never needs
to be differentiated through the Hamiltonian.
This decoupling is a direct computational benefit of the
port-Hamiltonian separation.

\subsection{Component-Wise Gradient Channels}
\label{subsec:ch4:grad_channels}

Each force channel contributes its own gradient pathway:
\begin{align}
  \frac{\partial\mathcal{L}}{\partial\beta}
  &= \sum_t\frac{\partial\mathcal{L}}{\partial F_{\mathrm{goal},t}}
     \cdot(-\nabla_{\vec{c}}D_\mathcal{G}),
  \label{eq:ch4:dL_dbeta} \\
  \frac{\partial\mathcal{L}}{\partial\alpha_j}
  &= \sum_t\frac{\partial\mathcal{L}}{\partial F_{\mathrm{bar},t}}
     \cdot(-\nabla_\Pos b_{\mathrm{IPC}}(d_j)),
  \label{eq:ch4:dL_dalpha} \\
  \frac{\partial\mathcal{L}}{\partial\lambda_s}
  &= \sum_t\frac{\partial\mathcal{L}}{\partial F_{\mathrm{soft},t}}
     \cdot(-\nabla_{\vec{c}}\tilde{r}),
  \label{eq:ch4:dL_dlambda_s} \\
  \frac{\partial\mathcal{L}}{\partial\lambda_h}
  &= \sum_t\frac{\partial\mathcal{L}}{\partial F_{\mathrm{hard},t}}
     \cdot(-\nabla_{\vec{c}}b_{\mathrm{sp}}(\phi)),
  \label{eq:ch4:dL_dlambda_h} \\
  \frac{\partial\mathcal{L}}{\partial\gamma}
  &= \sum_t\frac{\partial\mathcal{L}}{\partial(-R_{\theta,t})}
     \cdot(-\Mass^{-1}\Mom_t).
  \label{eq:ch4:dL_dgamma}
\end{align}

The Jacobians in each equation are \emph{analytical}:
$\nabla\tilde{r}$, $\nabla b_{\mathrm{IPC}}$, $\nabla b_{\mathrm{sp}}$
are known functions of the oracle patch.
No finite-difference approximation is needed.
This means each gradient channel carries interpretable structure:
$\partial\mathcal{L}/\partial\lambda_s$ is non-zero only where the
risk gradient $\nabla\tilde{r}$ is large; it points the coefficient
update toward configurations where risk matters.

% =======================================================================
\section{Double Regularization: Energy-Level and Field-Level}
\label{sec:ch4:regularization}
% =======================================================================

Learning to reshape a force field requires regularizing not just
the Hamiltonian values but the field's spatial structure.
Two levels are used.

\paragraph{Energy-level regularization $\mathcal{R}_H$.}
\begin{equation}
  \mathcal{R}_H(\theta)
  = w_\gamma|\gamma-\gamma_0|^2
  + w_\lambda(\|\lambda_s\|^2+\|\lambda_h\|^2).
  \label{eq:ch4:R_H}
\end{equation}
The first term anchors the dissipation near a reference value,
enforcing passivity (Definition~\ref{def:passivity}).
The second prevents material heads from dominating geometry
before the risk distribution has been sufficiently sampled.

\paragraph{Gradient-field regularization $\mathcal{R}_{\nabla H}$.}
\begin{equation}
  \mathcal{R}_{\nabla H}(\theta)
  = w_c\mathcal{L}_{\mathrm{clear}}
  + w_m\mathcal{L}_{\mathrm{multi}}.
  \label{eq:ch4:R_nablaH}
\end{equation}
$\mathcal{L}_{\mathrm{clear}}$ penalizes barrier violations along
rollouts; $\mathcal{L}_{\mathrm{multi}}$ over-samples near-contact
starts, sharpening the barrier gradient near $\{d_j=0\}$.
Together they enforce spatial smoothness of the force field near
constraint boundaries without them, large spatial gradients in
$\nabla\tilde{r}$ near terrain class boundaries cause the trajectory
to oscillate as it crosses material transitions.

The principled version is the Jacobian penalty
$\mathcal{R}_{\mathrm{Jac}} = w_J T^{-1}\sum_t
\|\nabla_\Pos F_{\mathrm{tot}}\|_F^2$;
the present implementation uses $\mathcal{R}_{\nabla H}$ as a
tractable surrogate.

\paragraph{Total loss.}
\begin{equation}
  \mathcal{L}_{\mathrm{total}}(\theta)
  = \mathcal{L}_{\mathrm{task}}
  + \mathcal{R}_H(\theta)
  + \mathcal{R}_{\nabla H}(\theta),
  \label{eq:ch4:L_total}
\end{equation}
where $\mathcal{L}_{\mathrm{task}} \in \{\mathcal{L}_{\mathrm{P1}},
\mathcal{L}_{\mathrm{P2}}\}$.

% =======================================================================
\section{Three Loops as Three Computational Jobs}
\label{sec:ch4:loops}
% =======================================================================

The three adaptation loops are not merely nested clocks.
They consume different information, produce different outputs, and
solve different optimization subproblems. Figure~\ref{fig:ch4:three_loops} summarizes the division of labor
across the three loops before the individual update rules are written
out in detail.

\begin{figure}[t]
    \centering
    \includegraphics[width=0.96\textwidth]{figs/ch4_three_loops.png}
    \caption{
    % \textbf{Three-loop update hierarchy as three distinct computational jobs.}
    The segment loop performs a local Gauss--Newton correction of the active coefficient vector using only segment-level observables and residual mismatch.
    The episode loop performs global optimization of the meta-navigator parameters $\xi$ using full-rollout statistics and the CVaR-based episode objective.
    The curriculum loop performs a statistical advancement decision using passivity and coverage records accumulated over many episodes.
    These loops are therefore not merely different timescales; they consume different information, update different objects, and solve different optimization subproblems.
    }
    \label{fig:ch4:three_loops}
\end{figure}

\subsection{Loop 1 -- Segment: Local Gauss-Newton Coefficient Correction}
\label{subsec:ch4:seg_loop}

\textbf{What it consumes.}
The observable residual from a single navigation segment:
$\Delta y_k^{\mathrm{des}} = y_k^\star - y_k$ where
$y_k = [-\mathrm{clr}_k, \mathrm{dist}_k, -\mathrm{speed}_k]^\top$.

\textbf{What it produces.}
An updated coefficient vector $\zeta_{k+1}$ for the \emph{next}
segment, via a Tikhonov-regularized Gauss-Newton step:
\begin{equation}
  \Delta\zeta_k = (J_k^\top J_k + \lambda_J I)^{-1}J_k^\top\Delta y_k^{\mathrm{des}},
  \qquad
  \zeta_{k+1} = \Pi_{\mathbb{R}_+}(\zeta_k + \kappa\Delta\zeta_k).
  \label{eq:ch4:seg_update}
\end{equation}
$J_k \in \mathbb{R}^{3\times|\zeta|}$ is the smoothed secant Jacobian,
updated with a rank-1 rule (same as Box~3.3 in Chapter~\ref{chap:grlsnam}).

\textbf{What it does not touch.}
The meta-navigator weights $\xi$ are unchanged.
The episode-level loss is not evaluated.
The conquest flags $\hat{c}_e$ are not updated.

\textbf{Why it cannot run at episode frequency.}
The residual $\Delta y_k^{\mathrm{des}}$ is a local signal from one
corridor segment.
Using it to update $\xi$ would overfit to the current segment geometry.
The Gauss-Newton solve is $(2+|\mathcal{C}_k|)$-dimensional,
not $|\xi|$-dimensional implying that it is explicitly \emph{not} a meta-update.

\subsection{Loop 2 -- Episode: Global CVaR Optimization}
\label{subsec:ch4:ep_loop}

\textbf{What it consumes.}
A full episode rollout: the trajectory $\{\PhaseState_t\}_{t=0}^T$
and the accumulated episode cost $J$ from \eqref{eq:ch4:J},
computed over all segments.

\textbf{What it produces.}
An update to the meta-navigator weights $\xi$ via the CVaR gradient:
\begin{equation}
  \xi \leftarrow \xi - \eta\nabla_\xi\mathcal{L}_{\mathrm{total}}.
  \label{eq:ch4:ep_update}
\end{equation}
The gradient $\nabla_\xi\mathcal{L}_{\mathrm{total}}$ is computed
through the rollout sensitivity recursion
\eqref{eq:ch4:dp_dtheta}--\eqref{eq:ch4:dq_dtheta}.

\textbf{What it does not touch.}
The curriculum sampling distribution $\Pi_{\mathrm{curr}}^e$ is
unchanged.
The phase flag (P1/P2) is unchanged.

\textbf{Why it cannot run at segment frequency.}
CVaR at level $\beta=0.95$ requires a \emph{batch} of episode costs
to estimate the 95th percentile; a single segment cost does not
define a tail distribution.
The rollout sensitivity recursion must propagate through the full
episode, so per-segment accumulation is necessary before the gradient
is formed.

\subsection{Loop 3 -- Curriculum: Statistical Conquest and Phase Transition}
\label{subsec:ch4:curr_loop}

\textbf{What it consumes.}
The per-episode passivity surrogate $\mathcal{L}_{\mathrm{pass}}$
and goal-distance records, accumulated over many episodes.

\textbf{What it produces.}
A conquest flag per start–goal pair:
\begin{equation}
  \hat{c}_e(\vec{q}_0,\vec{g})
  = \mathbf{1}\!\bigl[
      \mathcal{L}_{\mathrm{pass}} \leq \epsilon_{\mathrm{pass}}
      \;\wedge\;
      |\vec{q}_T - \vec{g}| \leq \epsilon_{\mathrm{goal}}
    \bigr],
  \label{eq:ch4:conquest}
\end{equation}
a coverage fraction
\begin{equation}
  \mathrm{cov}(\mathcal{D}_e) = |\mathcal{D}_e|^{-1}\sum\hat{c}_e,
  \label{eq:ch4:coverage}
\end{equation}
a re-weighted sampling distribution
$\Pi_{\mathrm{curr}}^{e+1}\propto 1-\hat{c}_e$,
and when advancement criterion \eqref{eq:ch4:advance} holds  the
phase transition P1→P2.

\textbf{What it does not touch.}
$\xi$ is unchanged within the curriculum update.
Individual coefficients $(\beta,\lambda_s,\lambda_h,\{\alpha_j\},\gamma)$
are unchanged.

\textbf{Why it cannot run at episode frequency.}
Conquest flags require enough episodes to assess \emph{stability}
(is $\mathcal{L}_{\mathrm{pass}}$ reliably small?) not just success
on one rollout.
The coverage fraction $\mathrm{cov}(\mathcal{D}_e)$ is a statistical
estimate that converges only over many episodes.
The phase transition is irreversible; triggering it on a single
episode's statistics would be unreliable.

\begin{equation}
  \textbf{Phase advance when:}\quad
  \mathcal{L}_{\mathrm{pass}} \leq \epsilon_{\mathrm{pass}}
  \;\text{ and }\;
  \mathrm{cov}(\mathcal{D}_e) \geq \rho_{\mathrm{cover}}.
  \label{eq:ch4:advance}
\end{equation}

% =======================================================================
\section{Factored Algorithms}
\label{sec:ch4:algorithms}
% =======================================================================

\begin{algorithm}[t]
\caption{Stage~2 Outer Curriculum Loop (with CVaR minibatch and moving-window transition)}
\label{algo:ch4:curriculum}
\begin{algorithmic}[1]
\STATE \textbf{Input:} trained Phase~1 weights $\xi_{\mathrm{P1}}$,
       oracle maps $Z^{(1)}, Z^{(2)}$, start--goal pool $\mathcal{D}$,
       batch size $B$, window size $W$, confidence level $\beta$
\STATE $\xi \leftarrow \xi_{\mathrm{P1}}$;\;
       $\lambda_s,\lambda_h \leftarrow 0$ (frozen);\;
       $\Pi^0_{\mathrm{curr}} \leftarrow \mathrm{Uniform}(\mathcal{D})$;\;
       $\mathrm{phase} \leftarrow \mathrm{P1}$;\;
       $\mathcal{B} \leftarrow \emptyset$ (episode cost buffer);\;
       $\mathcal{W} \leftarrow \emptyset$ (passivity window buffer)
\FOR{curriculum episode $e = 1,2,\ldots$}
  \STATE $(\vec{q}_0,\vec{g}) \sim \Pi^e_{\mathrm{curr}}$
  \STATE $(J_e,\mathcal{L}_{\mathrm{pass},e})
         \leftarrow \textsc{EpisodeProcedure}(\vec{q}_0,\vec{g},\xi,\mathrm{phase})$
         \hfill $\triangleright$ Algorithm~\ref{algo:ch4:episode}
  \STATE $\mathcal{B} \leftarrow \mathcal{B} \cup \{J_e\}$;\;
         $\mathcal{W} \leftarrow \mathcal{W} \cup \{\mathcal{L}_{\mathrm{pass},e}\}$
         (retain last $\max(B,W)$ entries)
  \IF{$|\mathcal{B}| \geq B$} \hfill $\triangleright$ Loop~2: CVaR batch update
    \STATE $\hat\nu \leftarrow \mathrm{VaR}_\beta(\mathcal{B})$ \hfill $\triangleright$ empirical $\beta$-quantile
    \STATE $\widehat{\mathrm{CVaR}}_\beta \leftarrow
           \hat\nu + \tfrac{1}{(1-\beta)|\mathcal{B}|}
           \sum_{J\in\mathcal{B}}(J - \hat\nu)_+$
    \STATE $\xi \leftarrow \xi - \eta\,\nabla_\xi\!\left[
           \widehat{\mathrm{CVaR}}_\beta + \mathcal{L}_{\mathrm{pass}} + \mathcal{L}_{\mathrm{reg}}
           \right]$ \hfill using indicator-weighted gradient
    \STATE $\mathcal{B} \leftarrow \emptyset$ \hfill reset buffer
  \ENDIF
  \STATE $\hat{c}_e(\vec{q}_0,\vec{g}) \leftarrow
         \mathbf{1}\!\left[
         \mathcal{L}_{\mathrm{pass},e}\leq\epsilon_{\mathrm{pass}}
         \wedge \|\vec{q}_T-\vec{g}\|\leq\epsilon_{\mathrm{goal}}
         \right]$
  \STATE Update $\Pi^{e+1}_{\mathrm{curr}} \propto 1-\hat{c}_e$
         \hfill $\triangleright$ Loop~3 sampling
  \IF{$\mathrm{phase}=\mathrm{P1}$ \textbf{and}
      $|\mathcal{W}| \geq W$ \textbf{and}
      $\frac{1}{W}\sum_{\mathcal{W}}\mathcal{L}_{\mathrm{pass}} \leq \epsilon_{\mathrm{pass}}$
      \textbf{and}
      $\mathrm{cov}(\mathcal{D}_e)\geq\rho_{\mathrm{cover}}$}
      \hfill $\triangleright$ moving-window passivity criterion
    \STATE $\mathrm{phase}\leftarrow\mathrm{P2}$;\;
           unfreeze $(\lambda_s,\lambda_h)$
           \hfill $\triangleright$ Loop~3 phase transition
    \STATE $\mathcal{W} \leftarrow \emptyset$ \hfill reset window
  \ENDIF
\ENDFOR
\end{algorithmic}
\end{algorithm}

\begin{algorithm}[t]
\caption{Episode Procedure (one forward rollout + Loop~2 gradient)}
\label{algo:ch4:episode}
\begin{algorithmic}[1]
\STATE \textbf{Input:} $(\vec{q}_0,\vec{g})$, weights $\xi$, phase, oracle $Z^{(1)},Z^{(2)}$
\STATE $\PhaseState_0 \leftarrow (\vec{q}_0,\mathbf{0})$;\;
       $J_{\mathrm{ep}} \leftarrow 0$;\;
       $\mathcal{L}_{\mathrm{pass}} \leftarrow 0$
\STATE Compute waypoints $\{g_k\}$ via geometry-only A$^*$ scaffold
\FOR{segment $k=1,\ldots,K$}
  \STATE Query $f_\theta(\mathcal{E}_k,P_k)\to
         (\beta,\{\alpha_j\},\lambda_s,\lambda_h,\gamma)$
         \hfill $\triangleright$ CoefEnergyNetMaterial forward pass
  \FOR{step $t=1,\ldots,T_k$}
    \STATE Assemble $F_{\mathrm{tot}}$ via
           \eqref{eq:ch4:F_goal}--\eqref{eq:ch4:R}
    \STATE $\Mom_{t+1}\leftarrow\Mom_t+\tau F_{\mathrm{tot}}$;\;
           $\vec{q}_{t+1}\leftarrow\vec{q}_t+\tau\Mass^{-1}\Mom_{t+1}$
    \STATE $J_{\mathrm{ep}} \leftarrow
           J_{\mathrm{ep}} + j_t(\vec{q}_t,\tilde{r},\phi)$
    \STATE $\mathcal{L}_{\mathrm{pass}} \leftarrow
           \mathcal{L}_{\mathrm{pass}} +
           \frac{\mathrm{softplus}(\Delta H_t)}{T}$
  \ENDFOR
  \STATE \textsc{SegmentUpdate}$(\PhaseState_k,y_k,J_{k-1})$
         \hfill $\triangleright$ Loop~1 / Procedure~\ref{proc:ch4:segment}
\ENDFOR
\STATE \textbf{Return:} $J_{\mathrm{ep}}, \mathcal{L}_{\mathrm{pass}}$
\end{algorithmic}
\end{algorithm}

\begin{mdframed}[linecolor=violet!60!black, linewidth=1pt,
                 frametitle={\small\sffamily\bfseries
                 Procedure 4.3: \textsc{SegmentUpdate} (Loop~1)}]
\label{proc:ch4:segment}
\textbf{Input:} $\PhaseState_k$, observable $y_k$, cached $J_{k-1}$\\
\textbf{Step 1:} Rank-1 Jacobian update (same as Box~3.3):
$J_k = \rho J_{k-1} + (1-\rho)\tilde{J}_k.$\\
\textbf{Step 2:} Gauss-Newton coefficient update \eqref{eq:ch4:seg_update}.\\
\textbf{Step 3:} If residual large, port correction (same as Box~3.3).\\
\textbf{Updates:} $\zeta_{k+1}$, cached $J_k$.\\
\textbf{Does not update:} $\xi$, phase flag, conquest flags.\\
\textbf{Cost:} $O(|\mathcal{C}_k|^3)$ (Tikhonov solve on active-set dimension).
\end{mdframed}

% =======================================================================
\section{Persistent State, Update Scope, and Efficiency}
\label{sec:ch4:efficiency}
% =======================================================================

\begin{table}[t]
\centering
\caption{%
  \textbf{Persistent state, update scope, and per-step cost
  in Stage~2.}
  Three categories: persistent across the full run, segment-local,
  and episode-local.
  Each row states what information triggers the update and what cost it incurs.
}
\label{tab:ch4:state}
\small
\setlength{\tabcolsep}{4pt}
\renewcommand{\arraystretch}{1.3}
\resizebox{\textwidth}{!}{
\begin{tabular}{lllll}
\toprule
\textbf{Object} & \textbf{Scope} & \textbf{Updated by} &
\textbf{Trigger} & \textbf{Cost} \\
\midrule
$\xi$ (meta-navigator) & Persistent & CVaR gradient (Loop~2) &
  After every episode & $O(|\xi|)$ \\
$\lambda_s,\lambda_h$ & Persistent & Phase transition (Loop~3) &
  Once at P1→P2 & $O(1)$ \\
Phase flag (P1/P2) & Persistent & Coverage criterion & Once & $O(1)$ \\
$\hat{c}_e$ (conquest flags) & Persistent & Per episode & Success + passivity &
  $O(|\mathcal{D}_e|)$ \\
$\Pi_{\mathrm{curr}}^e$ (sampling) & Persistent & Renorm from $\hat{c}_e$ &
  After every episode & $O(|\mathcal{D}_e|)$ \\
\midrule
$\zeta_k=(\eta_k,\mu_k)$ (coefficients) & Segment-local &
  Gauss-Newton (Loop~1) & After every segment & $O(|\mathcal{C}_k|^3)$ \\
$J_k$ (Jacobian) & Segment-local &
  Rank-1 secant & After every step & $O(3\,|\mathcal{C}_k|)$ \\
\midrule
$J_{\mathrm{ep}}$ (episode cost) & Episode-local &
  Accumulated per step & Each integration step & $O(1)$ \\
$\mathcal{L}_{\mathrm{pass}}$ (passivity) & Episode-local &
  softplus$(\Delta H_t)$ & Each integration step & $O(1)$ \\
\bottomrule
\end{tabular}}
\end{table}

\paragraph{Computational efficiency.}
The dominant per-step costs are the barrier gradient evaluations
$O(|\mathcal{C}_k|)$ and the rank-1 Jacobian update $O(3|\mathcal{C}_k|)$.
The Tikhonov solve is $(2+|\mathcal{C}_k|)$-dimensional
(typically 5--12 in practice), not $|\xi|$-dimensional.
The meta-navigator $f_\theta$ runs once per segment as a forward pass.
No per-step optimization occurs.

\paragraph{What is never recomputed from scratch.}
The Jacobian $J_k$ is always a rank-1 update of $J_{k-1}$.
The conquest flags $\hat{c}_e$ are maintained as a running record, they are not recomputed from the episode archive.
The curriculum sampling distribution is a simple renormalization
of $1-\hat{c}_e$.

% =======================================================================
\section{Oracle Instantiation}
\label{sec:ch4:implementation}
% =======================================================================

The realization of abstract coefficient predictor $f_\theta(\mathcal{E},P_t)$ extends the
geometry backbone of Chapter~\ref{chap:grlsnam} with a risk encoder
and two material heads:

\begin{itemize}
  \item \textbf{Geometry Transformer.}
        Operates on the token set $\{(d_j,\vec{n}_j)\}$ of active
        obstacles, producing geometry embedding $\vec{h}_{\mathrm{geom}}$.
        Outputs: $(\beta,\{\alpha_j\},\gamma)$.
  \item \textbf{Risk CNN encoder.}
        Operates on the 6-channel patch $P_t\in\mathbb{R}^{6\times P\times P}$,
        producing risk embedding $\vec{h}_{\mathrm{risk}}$.
  \item \textbf{Material heads.}
        Two linear layers from
        $[\vec{h}_{\mathrm{geom}};\vec{h}_{\mathrm{risk}}]$
        outputting $\lambda_s\geq 0$, $\lambda_h\geq 0$
        via softplus activations.
\end{itemize}

In Phase~1, $(\lambda_s,\lambda_h)$ are frozen at zero; the risk
encoder and material heads receive no gradient.
In Phase~2, they are unfrozen and trained via the CVaR gradient.
Phase~3 (replacing oracle $\tilde{r}$ with learned $\hat\rho$)
adds a perception backbone as a fourth branch: the enrichment
chain of §\ref{sec:ch5:limits} applied one step further.

% =======================================================================
\section{What the Richer Learner Achieves}
\label{sec:ch4:results}
% =======================================================================

We summarize here the principal empirical outcomes of the richer learner in
the oracle-material Stage~2 setting. The model is evaluated on the DFC2018
Houston aerial benchmark over 20 high-risk test episodes, with oracle access
to the soft material-risk field $\tilde r$ and hard-hazard boundary field
$\phi$. The goal of this section is not to provide a full experimental
analysis, but to state clearly what the enriched learner adds beyond the
geometry-only foundation.

The most important result is that the Stage~2 learner substantially reduces
risk exposure while preserving nearly the same path efficiency as the
Stage~1 geometry-only model. Relative to the Stage~1 learner, the enriched
model reduces cumulative risk exposure by \textbf{56.4\%}. Relative to a
material-blind Dijkstra baseline, it achieves a \textbf{17.4\%} reduction in
risk, even though both methods follow waypoint structure derived from the same
underlying geometry. This shows that the added material-aware energy terms do
not merely perturb the geometry backbone; they materially change the induced
vector field in a way that improves safety.

The hard-hazard barrier term also produces a clear qualitative benefit. In the
full Stage~2 model, the learned policy exhibits \textbf{0 hard-hazard entries},
whereas the ablation without $H_{\mathrm{mat}}^{\mathrm{hard}}$ incurs
\textbf{7} such entries over the same evaluation set. This confirms that the
hard-hazard component is not redundant with the soft material-risk term:
it serves as a distinct steep-repulsion mechanism for strongly discouraging
dangerous regions that should not be traversed at all.
Formal forward invariance would additionally require a discrete CBF
projection or a sufficiently small timestep; empirically, the barrier alone
reduces hard-hazard entries to zero on this evaluation set.

The richer learner also improves dynamical quality, not only safety. The number
of oscillatory corrections drops by approximately \textbf{$7\times$}, from
$8.7$ to $1.2$ per episode. This reduction is attributable to the joint effect
of the explicit dissipative port $R_\theta$ and the risk-sensitive curriculum,
which together produce smoother local corrections and less brittle behavior near
risk and obstacle boundaries. Consistent with this, the passivity criterion is
satisfied at convergence, with
\[
\mathcal L_{\mathrm{pass}} = 0.008 \le \epsilon_{\mathrm{pass}} = 0.01.
\]
Thus, the enriched learner does not gain safety by introducing unstable or
overreactive updates; it remains within the desired dissipative regime.

Importantly, these gains are obtained at very small geometric cost. The average
path length increases by only \textbf{1.1\%} relative to the Stage~1 learner,
indicating that the reduction in risk is achieved through modest local detours
rather than overly conservative rerouting. In this sense, the richer learner
improves the quality of the decision dynamics rather than simply trading safety
for inefficiency.

Taken together, these results validate the central claim of this chapter:
enlarging the learner with material-aware energy terms, tail-risk objectives,
and multi-timescale updates yields a safer and smoother policy than the
geometry-only foundation, while preserving the efficiency of the original
Hamiltonian navigation backbone.

Comprehensive ablations, comparative tables, and failure analysis are deferred
to Chapter~\ref{chap:experiments}.

% =======================================================================
\section{Chapter Summary}
\label{sec:ch4:summary}
% =======================================================================

This chapter enlarged the geometry-only learner along four coupled axes:

\begin{enumerate}
  \item \textbf{Stored energy.}
        $H_\theta$ gains $H_{\mathrm{mat}}^{\mathrm{soft}}$ and
        $H_{\mathrm{mat}}^{\mathrm{hard}}$; $R_\theta$ is made an
        explicit separate port (\S\ref{sec:ch4:energy}).
  \item \textbf{Coefficient space.}
        $(\lambda_s,\lambda_h)$ are added, each with its own gradient
        channel through the rollout sensitivity recursion
        (\S\ref{sec:ch4:gradients}).
  \item \textbf{Objective.}
        CVaR tail-risk replaces expected cost, directly penalizing
        the catastrophic tail of the episode-cost distribution
        (\S\ref{sec:ch4:objective}).
  \item \textbf{Adaptation hierarchy.}
        Three loops: segment Gauss-Newton, episode CVaR descent,
        curriculum conquest, each solving a different computational
        subproblem with different information and different outputs
        (\S\ref{sec:ch4:loops}).
\end{enumerate}

The three loops do not add an external evaluator.
They enrich the internal one: the costate $\Mom$ now accumulates
material-risk geometry through two new force channels, monitors a
richer energy balance, and receives gradient signals from a
tail-risk objective and a statistically-grounded curriculum.

Chapter~\ref{chap:hierarchy} will show that this four-axis
enlargement is an instance of a general enrichment principle
that can extend the framework to further stages.

% -----------------------------------------------------------------------
\subsection*{Supporting Results for Stage 2}
% -----------------------------------------------------------------------

\begin{theorem}[Passivity preservation under backward-compatible initialization]
\label{thm:ch4:passivity_init}
Suppose Stage~1 satisfies $\mathcal{L}_{\mathrm{pass}} \leq \epsilon$ at
convergence, and Stage~2 is initialized with $\lambda_s = \lambda_h = 0$,
the same $\xi$, the same integrator, and the same port input.
Then at initialization, $\mathcal{L}_{\mathrm{pass}}^{(2)} = \mathcal{L}_{\mathrm{pass}}^{(1)} \leq \epsilon$.
\end{theorem}

\begin{proof}
With $\lambda_s = \lambda_h = 0$, $H_\theta^{(2)} = H_\theta^{(1)}$,
so all rollout statistics are identical.
\end{proof}

\begin{theorem}[Small-enrichment passivity continuity]
\label{thm:ch4:passivity_cont}
Suppose Stage~1 satisfies $\mathcal{L}_{\mathrm{pass}} \leq \epsilon$ and
the new energy term has bounded gradient $\|\nabla H_{\mathrm{new}}\| \leq G$.
For $|\lambda_{\mathrm{new}}| \leq \delta$,
\[
  \mathcal{L}_{\mathrm{pass}}^{(k+1)}
  \;\leq\;
  \mathcal{L}_{\mathrm{pass}}^{(k)} + C\,\delta + O(\Delta t^2),
\]
where $C$ depends on $G$ and the trajectory length.
\end{theorem}

\begin{proposition}[PAC coverage certificate for phase transition]
\label{prop:ch4:coverage_pac}
If $m$ start--goal pairs are attempted and $s$ are conquered (passivity
and success criteria satisfied), then with probability $1-\delta$:
\[
  p_{\mathrm{conquer}} \;\geq\; \mathrm{LCB}_\delta(s, m),
\]
where $\mathrm{LCB}_\delta$ is the Clopper--Pearson lower confidence bound.
The phase transition from P1$\to$P2 is statistically grounded when
$\mathrm{LCB}_\delta(s,m) \geq \rho_{\mathrm{cover}}$ over a moving window
of $W$ episodes, preventing advancement on a single lucky episode.
\end{proposition}